\subsection*{类型系统补记：\texttt{mdspan} 与 \texttt{ituple}}

选择一个分块需要知道张量形状、各轴划分为多少块，以及当前块的坐标。Croqtile 用形状、元组和有界索引显式表示这些信息，让分块表达式保留原本分散在偏移计算中的关系。

带形状声明 \co@s32 [128, 256] lhs@ 同时确定元素类型和秩。这里的尺寸是常量，编译器能在执行前求值，拒绝用三个坐标访问二维张量，并检查分块形状的兼容性。声明使用符号维度时，即使部分尺寸来自调用方，秩信息仍然保留。

\co@span(i)@ 语法取出一个维度：\co@lhs.span(0)@ 是 128，\co@rhs.span(1)@ 是 256。简写 \co@lhs.span@ 则复制整个形状。

\co@ituple@ 将整数值组成元组。在 \co@parallel {px, py} by [8, 16]@ 中，坐标还带有范围：当前值标识某个实例，上界描述整个网格。编译器可以化简上界及派生表达式，但并不是给所有实例确定同一个当前索引。依赖实际执行的索引计算仍在运行时进行。

第 2.4 和 2.5 节介绍具体声明。这里需要区分程序本身已知的信息与某次调用提供的信息。编译时检查使用前者，入口检查和运行时断言则处理依赖后者的约束。
